% titel: Blatt 8, A9: CsCl-Kantenlänge und Abstand
% Teil II: der durchgerechnete Weg.

\subsection*{Gegeben und gesucht}
$\rho(\mathrm{CsCl}) = 3.99\ g\,cm^{-3}$. Gesucht: Kantenlänge $a$ der
Elementarzelle und kürzester Abstand Cs$^+$--Cl$^-$.

\subsection*{Ansatz}
Im CsCl-Gitter $Z=1$, $M = 132.9 + 35.45 = 168.36\ g\,mol^{-1}$.

\subsection*{Rechnung}
\[
  a = \left(\frac{Z\,M}{N_A\,\rho}\right)^{1/3}
    = \left(\frac{1\cdot168.36}{6.022\cdot10^{23}\cdot 3.99}\right)^{1/3}
    = \left(7.01\cdot10^{-23}\,cm^3\right)^{1/3}
    \approx 4.12\cdot10^{-8}\,cm .
\]
Umrechnung: $1\,cm = 10^{10}\,pm$, also
\[
  a \approx 412\ pm.
\]
Kürzester Abstand (auf der Raumdiagonale, halbe Diagonale):
\[
  d = \frac{a\sqrt{3}}{2}
    = \frac{412\cdot 1.732}{2} \approx 357\ pm.
\]

\begin{ergebnis}
$a \approx 4.12\cdot10^{-8}\,$cm $\approx 412\,$pm;
kürzester Cs$^+$--Cl$^-$-Abstand $\approx 357\,$pm $\approx 3.6\,$Å.
\end{ergebnis}

\begin{kontrolle}
Der gefundene Abstand (357 pm) liegt plausibel im Bereich typischer
Ionenbindungslängen. Mit $\rho\approx4\ g/cm^3$ und M$\approx168$ ergibt
sich ein ``sinnvolles`` Zellvolumen, keine negativen oder absurden Werte.
\end{kontrolle}

\begin{merksatz}
$\rho = ZM/(N_A a^3)$ nach $a$ auflösen. Im CsCl-Gitter berühren sich
Cs$^+$ und Cl$^-$ auf der halben Raumdiagonalen, nicht entlang der Kante.
\end{merksatz}
